\documentclass{article}
\usepackage[utf8]{inputenc}
\usepackage[T1]{fontenc}
\usepackage[margin=2cm]{geometry}
\usepackage{hyperref}
\sloppy
\begin{document}
\section*{Historical source}
Encyclopaedia Britannica scan, source row 383. Parse 3.5 extracted the text. Super 3.5 VL corrected it from an enlarged crop of each block. It can still contain recognition errors.
Page headers, footers, and page numbers are removed.

\subsection*{POLITICAL ECONOMY (E0020)}
POLITICAL ECONOMY.

Chap. II.

The second fort is of those which human industry can multiply in proportion to the demand. Where the commodity, as corn, is such as cannot be produced but by human industry, the price is more uniform than in almost any other case. The increasing scarcity and consequently value of land, tends indeed to raise it; but this is counteracted by the invention of machinery, and improved methods of labour. The opposite agency of these two causes has a constant tendency to preserve uniformity in the value of grain; though we cannot, with Smith, consider this uniformity as likely to be so complete, as to render the price of grain a sure standard for the value of silver.

There are other commodities, however, which nature produces in abundance, or which, where land is plentiful, can be multiplied with little or no cultivation. Of these the principal is butcher meat. Lands can be covered with cattle or sheep by the labour of few hands, and sometimes without any labour at all. Hence, in rude times, butcher meat is always cheaper than corn; in improved periods, the reverse is the case. For a long time the price continues constantly to rise, as we have seen it do throughout Great Britain, the pasture lands being more and more converted into arable. At last, however, it becomes so high as to make it an object for the farmer to fall his cattle, and to cultivate ground for the purpose of feeding them. After this era, the price is likely to experience a certain diminution, from the improved modes of feeding and rearing, which, in consequence of this new attention, are likely to be discovered and adopted.

covered and adopted.
There are certain animals, as hogs, poultry, \&c.
which are fed on mere offals, and in a rude state, there-
fore, are full cheaper than butcher meat. In an im-
proved state they are dearer; for they have not as yet,
at least in this country, become an object of separate
cultivation.

The third sort consists of those, in the multiplication of which the power of man is either limited or uncertain. In these the rule is various. Some commodities are not cultivated on their own account, but are appendages to others; as wool and hides to the carcase of the ox or fleece. Both these commodities are much more portable, and more easily preserved, than the flesh of the animals from which they are taken; the market for them is thus much more extensive, and the demand more equal at all times. Hence, in rude periods, when the flesh of animals, from its abundance, is of small value, the appendages equal or surpass it in price. At Buenos Ayres frequently, and sometimes even in Spain, an ox is killed for the sake of the hide and tallow. In an improved state of society, on the contrary, the hide and fleece become considerably inferior in value to the carcase.

Fish is an article, the supply of which is considerably limited, as man has no power of production in respect to it, though, by the exertion of industry, he can collect a greater quantity. Shoals of fish are generally copious, but uncertain.

pious, but uncertain.
Metals and minerals are articles, the supply of which
is not precisely limited, but extremely uncertain. The
discovery of new mines, or the continuance of fertility
in the old, are equally beyond the reach of calcula-
tion.

CHAP. II. Of the Sources of Wealth.

All wealth arises from three sources; it is either produced by the spontaneous bounty of nature, or it is the fruit of human industry, or it is generated by the judicious employment of a quantity of wealth previously accumulated. To these three heads then of land, labour, and capital, all national wealth may be referred.

Smith has treated of the revenue derived from these three sources as forming the constituent parts of the price of commodities; and with regard to labour in particular, repeatedly considers it as the only source of wealth*. According to the view, however, given above, the price of all commodities depends entirely on the proportion between the demand and the supply. Labour, therefore, (and the same may be said of land and capital), is only a means of furnishing or increasing a supply of those articles for which there already exists a demand, and unless it be successful in so doing, the most severe labours will meet with no remuneration whatever. We shall therefore proceed to consider the revenue which arises from these different sources, and the circumstances by which it is increased or diminished.

Sect. I. Land.

All land which is not naturally barren, and is cultivated with any crop, affords something more than is necessary to pay the expense of labouring it. This surplus goes as a rent to the landlord, who, in consideration of receiving it without risk or trouble, relinquishes to the farmer the profits of cultivation.

The proportion of the produce of a field which is to go for rent, varies with different circumstances. The chief of these is the fertility of the soil, the extent of the market, which enables the produce to be disposed of to greater advantage, the propriety or poverty of the country, which causes a greater or less demand for that produce, and the average skill and activity of the farmers, which will enable them to turn the fertility of the ground to better account. It is almost needless to observe, were it not for the vague language often made use of upon this subject, that the rate at which farms let, must, like all other commodities, depend altogether upon the demand and the supply. If much is to be made by farming, many will bid for farms, and the rent will be raised by their competition, and vice versa. The idea that all the landlords of an extensive country may combine to raise their rents, is altogether chimerical. Even could it take place, it could be accomplished only by a certain number of them allowing their lands to lie waste, which, diminishing the supply, would doubtless raise the rent of the cultivated lands. But we need not fear that any landlord should leave his lands in this condition, from a culpable scheme of aggrandizing the rest of his body at his own expense, as well as that of the public.

Land which produces food for man will at all times afford rent to the landlord, in proportion to its fertility, and the other circumstances mentioned above. Men multiply in proportion to the means of subsistence; they have even a constant tendency to multiply beyond those means; hence there is always a full demand for this species


\section*{Source record}
\begin{description}
\item[Scan ID] \texttt{britannica:309f765b1ecc83e0a7f7762e44b4a927a7555443:train:383}
\item[Image file] \texttt{page-0383-018fb71412b2.jpg}
\item[Image SHA-256] \texttt{018fb71412b2c0c66f9e9afefa55448a8bc9910b962f68d60bceb3b17b3b8161}
\item[Dataset] NationalLibraryOfScotland/encyclopaedia\_britannica\_illustrated at revision \texttt{309f765b1ecc83e0a7f7762e44b4a927a7555443} (CC0-1.0)
\item[Extraction] model \texttt{nvidia/nemotron-3.5-parse}, raw reply \texttt{extract/raw/page-0383-018fb71412b2.json}
\item[Raw reply SHA-256] \texttt{de0ea22dade366596d957dd43057a0a477c2abadfc582ef8eb20ea259aa0f6dd}
\item[Correction] model \texttt{nvidia/Nemotron-3.5-Super-VL}, block crops, raw replies \texttt{correct/raw/crop\_nothink/page-0383-018fb71412b2-b*.json}; super\_vl: 15
\item[Running header] 110 Nature of Wash, Sc.
\end{description}
\end{document}
